
The dual diagram \hyperref[fig:sga3-exp5-dual-012star]{\((0,1,2)^\ast\)} considered at the beginning of this
proof induces the following commutative diagram:

\SGAThreeDiagram[0.78\textwidth]
  {figures/exp5/p267_proof_iv_ring_modules_diagram.png}
  {The commutative ring-and-module diagram concluding the affine proof
  of \hyperref[thm:sga3-V-4.1-iv]{Theorem~4.1(iv)}.}
  {fig:sga3-exp5-p267-proof-iv-ring-modules}
  {LNM 151, p.~267}

Here \(u_0,u_1,u_2\) have, respectively, the component pairs
\[
  (i,\varepsilon),\qquad
  (\delta_1,\delta_0\varepsilon),\qquad
  (\delta'_2,\delta'_0\delta_0\varepsilon).
\]
We know that \(u_1\) is an isomorphism. Since the two left-hand
squares of \((0,1,2)^\ast\) are cocartesian, \(u_2\) is bijective.
The two horizontal rows of the diagram are exact, and therefore
\(u_0\) is bijective.%
\footnote{\label{note:sga3-V-26}Editor's note: the argument which follows has been added.}
Thus \(A_0\) is a locally free \(B\)-module of rank \(n\). By the
preceding reductions, it follows that
\[
  \delta_0\otimes\delta_1:
  A_0\otimes_BA_0\xrightarrow{\ \sim\ }A_1
\]
is an isomorphism. This completes the proof of \hyperref[thm:sga3-V-4.1]{Theorem~4.1} in the
special case under consideration: \(X_0\) is affine and \(d_1\) is
locally free of constant rank \(n\).

\section*{5. Passage to the quotient by a finite flat groupoid
  (the general case)}
\phantomsection\label{sec:sga3-V-5}
\addcontentsline{toc}{section}{5. Passage to the quotient by a finite
  flat groupoid (the general case)}

\phantomsection\label{par:sga3-V-5-a}%
\noindent\textbf{a)}
Let \(U^{(n)}\) be the largest open subset of \(X_0\) over which
\(d_1\) is finite locally free of rank \(n\). The scheme \(X_0\) is
the disjoint union of the \(U^{(n)}\). The two Cartesian squares

\begin{figure}[H]
  \centering
  \resizebox{0.29\textwidth}{!}{%
    \SGAThreeNativeInput{figures/exp5/p267_section5_rank_square_d0.png}}
  \qquad\textit{and}\qquad
  \resizebox{0.29\textwidth}{!}{%
    \SGAThreeNativeInput{figures/exp5/p267_section5_rank_square_d1.png}}
  \caption{The two Cartesian squares controlling the rank strata
    \(U^{(n)}\).}
  \label{fig:sga3-exp5-p267-section5-rank-squares}
  \par\smallskip
  {\footnotesize\itshape Source: LNM 151, p.~267.}
\end{figure}

\noindent show that the inverse images of \(U^{(n)}\) under \(d_0\)
and \(d_1\) both coincide with the largest open subset of \(X_1\)
over which \(d'_2\) is locally free of rank \(n\).%
Thus
\[
  d_0^{-1}(U^{(n)})=d_1^{-1}(U^{(n)}),
\]
and the groupoid \(X_\ast\) is the disjoint union of the groupoids
\(X_\ast^{(n)}\) induced by \(X_\ast\) on the open-and-closed
subschemes \(U^{(n)}\). It is therefore enough to prove
\hyperref[thm:sga3-V-4.1]{Theorem~4.1} for each \(X_\ast^{(n)}\). We are reduced to the case in
which \(d_1\) is finite locally free of rank \(n\).

\medskip
\phantomsection\label{par:sga3-V-5-b}%
\noindent\textbf{b)}
\textit{We can now prove \hyperref[thm:sga3-V-4.1]{the theorem} in the general case.}
